% from ega2-1-fr.tex
\oldpage[II]{72}

\begin{proposition}[4.1.4]
\label{II.4.1.4-fr}
Soit $\mathcal{L}$ un $\mathcal{O}_Y$-Module inversible. Pour tout
$\mathcal{O}_Y$-Module quasi-cohérent $\mathcal{E}$, il existe un
$Y$-isomorphisme canonique
$i_{\mathcal{L}}:\mathbf{P}(\mathcal{E})\xrightarrow{\sim}
\mathbf{P}(\mathcal{E}\otimes\mathcal{L})$~; en outre, si l'on pose
$P=\mathbf{P}(\mathcal{E})$, $Q=\mathbf{P}(\mathcal{E}\otimes\mathcal{L})$,
$i_{\mathcal{L}}^*(\mathcal{O}_Q(n))$ est canoniquement isomorphe à
$\mathcal{O}_P(n)\otimes_Y\mathcal{L}^{\otimes n}$ pour tout
$n\in\mathbf{Z}$.
\end{proposition}

\begin{proof}
Remarquons d'abord que si $A$ est un anneau, $E$ un $A$-module, $L$ un
$A$-module \emph{monogène libre}, on définit canoniquement un homomorphisme
de $A$-modules
\[
  \mathbf{S}_n(E\otimes L)\longrightarrow
  \mathbf{S}_n(E)\otimes L^{\otimes n}
\]
en faisant correspondre à
$(x_1\otimes y_1)\cdots(x_n\otimes y_n)$ l'élément
\[
  (x_1x_2\cdots x_n)\otimes(y_1\otimes y_2\otimes\cdots\otimes y_n)
  \qquad(x_i\in E,\ y_i\in L\text{ pour }1\leq i\leq n)~;
\]
on vérifie immédiatement (en se ramenant au cas où $L=A$) que cet
homomorphisme est en fait un isomorphisme. On en conclut un isomorphisme
canonique de $A$-algèbres graduées
$\mathbf{S}_A(E\otimes L)\xrightarrow{\sim}
\bigoplus_{n\geq0}\mathbf{S}_n(E)\otimes L^{\otimes n}$. En revenant aux
conditions de (4.1.4), les remarques précédentes permettent de définir un
isomorphisme canonique de $\mathcal{O}_Y$-Algèbres graduées
\[
\label{II.4.1.4.1-fr}
  \mathbf{S}_{\mathcal{O}_Y}(\mathcal{E}\otimes_{\mathcal{O}_Y}\mathcal{L})
  \xrightarrow{\sim}
  \bigoplus_{n\geq0}\mathbf{S}_n(\mathcal{E})
  \otimes_{\mathcal{O}_Y}\mathcal{L}^{\otimes n}
\tag{4.1.4.1}
\]
en définissant cet isomorphisme comme un isomorphisme de préfaisceaux et
tenant compte de (\hyperref[II.1.7.4-fr]{1.7.4}),
(\hyperref[I.1.3.9-fr]{I, 1.3.9}) et
(\hyperref[I.1.3.12-fr]{I, 1.3.12}). La proposition résulte alors de
(\hyperref[II.3.1.8-fr]{3.1.8, (iii)}) et
(\hyperref[II.3.2.10-fr]{3.2.10}).
\end{proof}

\begin{env}[4.1.5]
\label{II.4.1.5-fr}
Avec les hypothèses de (\hyperref[II.4.1.1-fr]{4.1.1}), posons
$P=\mathbf{P}(\mathcal{E})$, et désignons par $p$ le morphisme structural
$P\to Y$. Comme on a par définition
$\mathcal{E}=(\mathbf{S}_{\mathcal{O}_Y}(\mathcal{E}))_1$, on a un
homomorphisme canonique
$\alpha_1:\mathcal{E}\to p_*(\mathcal{O}_P(1))$
(\hyperref[II.3.3.2.2-fr]{3.3.2.2}) et par suite aussi
(\hyperref[0.4.4.3-fr]{0, 4.4.3}) un homomorphisme canonique
\[
\label{II.4.1.5.1-fr}
  \alpha_1^\sharp:p^*(\mathcal{E})\longrightarrow\mathcal{O}_P(1).
\tag{4.1.5.1}
\]
\end{env}

\begin{proposition}[4.1.6]
\label{II.4.1.6-fr}
L'homomorphisme canonique (\hyperref[II.4.1.5.1-fr]{4.1.5.1}) est surjectif.
\end{proposition}

\begin{proof}
En effet, on a vu dans (\hyperref[II.3.3.2-fr]{3.3.2}) que
$\alpha_1^\sharp$ correspond fonctoriellement à l'homomorphisme canonique
\[
  \mathcal{E}\otimes_{\mathcal{O}_Y}
  \mathbf{S}_{\mathcal{O}_Y}(\mathcal{E})
  \longrightarrow
  (\mathbf{S}_{\mathcal{O}_Y}(\mathcal{E}))(1)~;
\]
comme par définition $\mathcal{E}$ engendre
$\mathbf{S}_{\mathcal{O}_Y}(\mathcal{E})$, cet homomorphisme est surjectif,
d'où la conclusion en vertu de (\hyperref[II.3.2.4-fr]{3.2.4}).
\end{proof}

\subsection{Morphismes d'un préschéma dans un fibré projectif.}
\label{subsection:II.4.2-fr}

\begin{env}[4.2.1]
\label{II.4.2.1-fr}
Gardant les notations de (\hyperref[II.4.1.5-fr]{4.1.5}), soient $X$ un
$Y$-préschéma, $q:X\to Y$ le morphisme structural, et soit $r:X\to P$ un
$Y$-\emph{morphisme}, de sorte que l'on a le diagramme commutatif
\[
  \xymatrix{
    P \ar[d]_p & X \ar[l]_r \ar[dl]^q
    \\Y
  }
\]

